## 5. Results — current findings

Numbers only. The plan/rules that produced these live in the methodology chapter (§4); this states what
happened when that plan was run. Cross-referenced from methodology by section number below.

### 5.1 Predictive models

- **Anchor-leak magnitude** (methodology §4.2): pooled vs. split-matched CR anchor differed by 3-11% (e.g. 4survey temporal: pooled 51.96m vs split-matched 46.48m top height).
- **`freeze_y_max` ablation** (§4.2): learned-`k` range shrinks to ~77% of the full two-parameter model's range when `y_max` is frozen — weakens but does not eliminate the "two knobs fighting" explanation.
- **Experiment A — physics-weight ablation** (§4.4): full weight (1.0) hurts in all 4 cohort×split combinations; `trajectory_weight ≥ 0.1` clearly hurts 4survey; nominal validation winner (`pw=0.1, tw=0.0`) does not beat the `0/0` control on held-out test (4survey: 0.6319 vs 0.6337; 6survey: 0.7468 vs 0.7483). **Verdict: no reliable predictive gain from physics.**
- **Experiment B — env-conditioned sweeps** (§4.4): E6 (90 fits) — `stage4_all_environmental` does not clearly beat smaller tiers. E7 (36 jobs, single split) — dropout shows no clear benefit. E9 (five-fold confirmation) — 4survey E7 winners hold within noise; 6survey winners were optimistic by 0.03-0.05 R² on single-split E7. Accuracy: PINN variants roughly on par with, or slightly behind, plain DNN.
- **`y_max`/`k` sign-flip** (§4.4): `pinn_env_terrain_k`'s learned `y_max`/`k` correlation flips sign between cohorts (-0.5748 on 4survey, +0.2096 on 6survey). **Verdict: evidence physics conditioning changes model behaviour, cohort-dependently — not proof of a clean, desirable effect.**
- **Neighbour-height leakage evidence** (§4.5): removing `neighbour_mean_height`/`neighbour_height_differential` dropped XGBoost spatial_block test R² from 0.598→0.321 (4survey) and 0.327→−0.337 (6survey) — confirms the leak was large, not theoretical.
- **Buffer-adequacy check** (§4.6): CR residual's own semivariogram range is ≈3,956-4,284 m — roughly two orders of magnitude larger than the 60 m training buffer.

### 5.2 Avenue 1 (`mean_cr_residual`)

- **NLME** (§4.7): 18 VIF-screened terrain/wind fixed effects, fit on `spatial_block` train. 17/18 individually significant (p<0.05; `tpi_500m` not, p=0.549). Largest standardised coefficients: `slope_degrees` +1.30, `topex` +0.98, `local_relief_500m` −0.96, `windward_topex` +0.70, `solar_radiation_index` +0.64. Compartment random-intercept variance explained: **2.02%** — real and broad across variables, not concentrated in one dominant driver.
- **Scope comparison** (§4.7, 2026-08-04, pooled five-fold): 4survey `terrain_and_wind_only` EN/XGB = 0.216/0.274, `all_environmental` EN/XGB = 0.360/0.411; 6survey near-zero/inconclusive for both scopes.
- **Paired fold-level significance** (§4.7, computed 2026-08-09): 4survey, `all_environmental` wins in 5 of 5 folds for both Elastic Net (mean Δ=+0.145, paired-t p=0.0005) and XGBoost (mean Δ=+0.139, paired-t p=0.0022) — Wilcoxon floored at p=0.0625 by n=5, but full directional agreement plus a consistent paired-t result is real, converging evidence. 6survey: direction inconsistent across folds for Elastic Net, not significant for either model. Full per-fold values: `outputs/spatial_block_kfold/cr_residual_environmental_paired_significance_folds.csv` / `_summary.csv`.
- **Spatial autocorrelation** (§4.7): CR residual Moran's I ≈0.1223, range ≈4,284 m — highest among all six models compared (i.e. the pooled-curve target leaves the most unexplained spatial structure of any model tested).

### 5.3 Avenue 2 (`local_y_max_difference`)

- **Representation checks** (§4.8): local shelter terms (`topex, windward_topex, whcl`) beat GWA wind-speed/Weibull variants; `tpi_500m` beat other terrain scales.
- **GNNWR evaluation** (§4.8, full training population, 5-fold pooled): R² ≈0.145 (`terrain_wind`), 0.318 (`terrain_wind_plus_management`), 0.286 (`broad_environment_plus_management`) — highest point estimate of the three model families tested, per scope. Earlier capped-reference-set runs (6,000-16,000 of 31,117 rows, GPU-memory workaround) were at best roughly tied with EN/XGBoost — superseded by the full-population numbers above.
- **Broad-environmental notebook** (§4.8, stale — not final): pre-fix CSVs (2026-08-04, before 19:04) reported `terrain_wind` = 0.125/0.117 (EN/XGB); fresh, post-fix reruns give 0.132/0.119 instead. Notebook itself not yet regenerated, so the full four-scope table is not final — only the qualitative claim (management added real signal, broader environment added comparatively little) is currently defensible.
- **Residual spatial autocorrelation** (§4.9): stays high (range 0.66-0.74, "~0.70") even for the best model (GNNWR) — real spatial structure left uncaptured by every model tested on this target.
- **Significance** (§4.9): paired cluster-bootstrap 95% CI on GNNWR's R² difference vs. EN/XGBoost crosses zero in all 4 pairwise comparisons tested — **GNNWR's lead is not statistically confirmed.**

### 5.4 Stale or superseded findings — do not cite as current

- Early DNN/PINN batch-size mismatch (128 vs 512) — superseded by the matched 256/256 rebuild (2026-07-30).
- Earlier 6-variable NLME model: 5.02% variance explained, included `elevation`, uncorrected REML variance-comparison — superseded by the 18-variable VIF-screened rebuild (2.02%, §5.2 above).
- E7 single-split 6survey R² values — optimistic by 0.03-0.05 vs. E9's five-fold confirmation (§5.1 above).
- Pre-fix Avenue 2 broad-environmental CSVs (2026-08-04, before 19:04) — `terrain_wind` = 0.125/0.117, superseded by fresh reruns (0.132/0.119); full table still not regenerated (§5.3 above).
- Earlier capped-reference-set GNNWR runs (6,000-16,000 rows) — "at best roughly tied" with EN/XGBoost, superseded by the full-population run (§5.3 above).
- `temporal_narrow_gap` numeric results from the retired `Top_Height99`+`yldc` pipeline (2026-07-13 to 2026-07-20) — not comparable to the current pipeline.
